\documentclass[11pt,a4paper]{article}
\usepackage[T1]{fontenc}
\usepackage[utf8]{inputenc}
\usepackage[french]{babel}
\usepackage{lmodern,microtype}
\usepackage[margin=23mm]{geometry}
\usepackage{booktabs,array,tabularx,enumitem}
\usepackage{hyperref}
\hypersetup{colorlinks=true,linkcolor=black,citecolor=black,urlcolor=black,
pdftitle={Raisonnement sans LLM pour l'audit IoT : références simples et modèles explicites},
pdfauthor={Projet LANCE - recherche assistée par Codex},pdflang={fr-FR}}
\urlstyle{same}
\def\UrlBreaks{\do\/\do\-\do\_\do\.\do\:}
\newcolumntype{Y}{>{\raggedright\arraybackslash}X}
\newcolumntype{P}[1]{>{\raggedright\arraybackslash}p{#1}}
\newcommand{\code}[1]{\nolinkurl{#1}}
\newcommand{\revision}{2026-10-07 14:28 (Europe/Brussels, UTC+02:00)}
\setlength{\parskip}{3pt}
\setlength{\parindent}{1.1em}
\setlength{\emergencystretch}{2em}
\setlist{nosep,leftmargin=5mm,topsep=3pt}
\clubpenalty=10000 \widowpenalty=10000
\setcounter{tocdepth}{1}
\title{Raisonnement sans LLM pour l'audit IoT\\
\large Références simples, politiques et vérification de configurations}
\author{Projet LANCE\\\small Recherche exploratoire assistée par Codex}
\date{\small Dernière révision : \revision\\[3pt]
Statut : revue ciblée et expériences proposées ; aucun gain mesuré}
\begin{document}
\maketitle
\begin{abstract}
Le comité de thèse demande de mesurer l'apport des modèles de langage face à
des méthodes plus simples. Cette revue examine le raisonnement défensif sur
dossiers et configurations fournis, sans nouvelle collecte réseau. Dix sources
primaires permettent de distinguer règles métier et requêtes relationnelles,
inférence Datalog, politiques d'autorisation, puis vérification symbolique de
configurations réseau. Le D1 à règles déjà documenté dans LANCE constitue le
point d'appui ; ce rapport propose d'en adapter les règles à une lecture hors
ligne de dossiers, avec normalisation commune, justification et abstention
explicite. Une base relationnelle facilite certaines liaisons ;
elle ne constitue pas à elle seule un nouveau raisonnement. Les formalismes
plus riches sont des options conditionnées par les propriétés étudiées et les
données disponibles. Quatre expériences hors ligne isolent moteur, formalisation
et information. Le coût humain de préparer les règles et les modèles doit être
observé. Aucun article n'a été reproduit et aucun résultat ne démontre encore
un avantage ou une équivalence de LANCE.
\end{abstract}
\tableofcontents
\bigskip
\noindent\textbf{Décision proposée.} Valider le D1 existant avant son adaptation
hors ligne. Conserver D0 secondaire et justifier chaque ajout par une propriété précise.

\section{Objectif et méthode de revue}
\textbf{Question.} Parmi les décisions d'interprétation et de vérification
documentaire de LANCE, lesquelles peuvent être obtenues par des règles simples,
à quelles limites, et quel comparateur permettrait ensuite d'attribuer un effet
au LLM ? Le périmètre inclut dossiers figés, exports de configuration, inventaires
et politiques fournis par le propriétaire, ainsi que des données synthétiques
signalées. Les expériences proposées évaluent des propriétés et des conclusions
hors ligne ; elles ne définissent aucun actionneur d'exploitation.

\textbf{Recherche ciblée.} Les requêtes ont porté sur SQL récursif, négation
Datalog, provenance Soufflé, OPA/Rego, Cedar, FIREMAN, Header Space Analysis,
Batfish et NetPlumber. Les documents ont été consultés le 7 octobre 2026.
Les passages exploités dépassent les résumés : sémantique et limites des langages,
sections de modélisation, protocoles d'évaluation et discussions des articles.
Les références précisent ces périmètres. Les sections non mentionnées ne sont
pas réputées avoir été relues intégralement. Les outils documentés ne sont pas
installés ni exécutés. Les temps et corpus différents des publications ne sont
pas mis en concurrence numérique.

\textbf{Documents locaux.} Lecture des instructions et index du dépôt, du
protocole agent et automatisation du 27 septembre et des rapports de méthodes du
7 octobre. Le protocole existant définit D0/D1, distingue attribution et utilité
pratique, et propose déjà un contraste sur observations communes. La présente
revue approfondit ses références classiques ; elle ne remplace pas ses mesures
de campagne. Le \href{../../docs/benchmark/agent-vs-automation-cli.md}{guide CLI}
documente un D1 déjà implémenté pour les phases 3 et 4,
analyse et vérification. Le \emph{lecteur à règles} décrit ci-dessous est un
adaptateur hors ligne \emph{proposé} de cette référence ; il ne décrit pas une
nouvelle capacité déjà livrée. Cette distinction est documentaire : aucun
exécutable D1 n'est validé par la présente revue. La base Git est
\code{a8dbeb9404f0cad960d61e42f02a1d70ca143a59}, avec modifications et conflits
locaux. Les empreintes de lecture sont consignées dans
\code{sans-llm-raisonnement-sources.json}. Le commit seul ne décrit pas cette
copie de travail ; aucun comportement exécuté n'est déduit des fichiers en conflit.

\section{Contrat : ce que les données permettent d'affirmer}
Une déclaration doit identifier la ressource, la propriété, l'identité, le
contexte d'autorisation et la période. Pour les propositions suivantes, une
observation conserve origine, date, version du producteur, périmètre et
transformations. Une identité non documentée ne devient pas « anonyme » et une
politique absente ne devient pas « accès interdit ». Une configuration décrivant
un droit attendu et une réponse observée décrivent deux objets différents.

\textbf{Politique proposée de complétude.} Chaque dossier déclare les champs
indispensables et ce qui est couvert : composants, attributs, période et
sous-ensemble de configuration. Une recherche négative ne permet une conclusion
que dans un domaine déclaré complet pour cette recherche. Sinon le verdict est
\emph{insuffisant}, avec obligation manquante. Il faut conserver des observations
positives et négatives explicites. La valeur \code{NULL} d'une base, l'absence
d'un tuple ou un résultat vide ne garantissent pas cette politique.

Les issues proposées sont \emph{étayé}, \emph{réfuté}, \emph{insuffisant},
avec \emph{conflit} signalé lorsqu'une contradiction persiste. Il reste à
préenregistrer la règle de priorité si le format impose un seul verdict.
Conserver la contradiction comme donnée est préférable à la perdre par
écrasement. Le moteur peut raisonner exactement sur un modèle incorrect ;
la correction formelle du calcul ne prouve donc pas la fidélité au terrain.

\section{Constats : règles, relations et inférence}
\subsection{Une référence forte peut rester ordinaire}
\textbf{Lecteur à règles proposé.} L'entrée commune est un dossier normalisé ; des règles
explicites appliquent obligations de preuve, compatibilité de contexte et
complétude. Le résultat cite les artefacts utilisés, les règles déclenchées,
les contradictions et les champs manquants. La politique peut tenir dans un
module Python ou des tables déclaratives. La lisibilité, la couverture et la
maintenance sont des mesures à établir ; le choix du langage ne les garantit pas.

D0 reste utile pour établir le niveau d'une procédure fixe correctement conçue,
mais D1 est le contraste principal du protocole existant : état, conditions
et abstention sont déjà
possibles sans LLM. Dans cette expérience hors ligne, « adaptation » signifie
actualiser un verdict ou ses obligations à partir du dossier reçu, sans acquérir
de nouvelle observation. Un gain contre D0 seul ne permet pas d'affirmer que
les méthodes conditionnelles classiques ont été dépassées.

\subsection{SQL : beaucoup de liaisons sont déjà des jointures}
\textbf{Fait sourcé.} SQLite fournit jointures, filtrage, agrégation et
déduplication. Un filtre \code{WHERE} élimine aussi les expressions évaluées
à \code{NULL} ; une jointure externe peut donc perdre ses lignes incomplètes
si le filtre est mal placé. Les expressions de table communes récursives
permettent également de parcourir des graphes ; l'exemple officiel utilise
\code{UNION} pour éviter une récursion cyclique répétitive~\cite{sql}.

\textbf{Proposition.} Relier une observation à la bonne déclaration via
dossier, ressource, identité et période ; dédupliquer ; lister les obligations
non satisfaites ; repérer deux observations incompatibles. Ces fonctions peuvent
être des jointures et des règles, sans base de graphes ni LLM. Les noms des tables
et identifiants ne doivent pas contenir les labels cachés. Les citations doivent
remonter aux observations et pas seulement à une ligne agrégée. SQL ne résout
pas automatiquement l'ambiguïté d'une identité ou l'intention métier.

Une variante relationnelle qui ajoute des champs absents du lecteur à règles améliore
l'information disponible. Celle qui ajoute une obligation améliore la politique.
Pour isoler la représentation, ce lecteur doit pouvoir consulter ces mêmes champs et
appliquer ces mêmes obligations. Si les verdicts restent identiques, une réduction
du temps de revue ou des erreurs d'attribution peut néanmoins être utile.

\subsection{Datalog/Soufflé : justifier la récursion et sa provenance}
\textbf{Faits sourcés.} Soufflé exprime relations typées et règles, dont une
récursion vers un point fixe et une négation stratifiée~\cite{souffle}.
L'article de Zhao et collaborateurs explique des arbres de dérivation construits
à partir d'annotations de règle et de hauteur. Il examine aussi les explications
de non-dérivation ; la négation signifie absence de tuple dans l'instance.
L'évaluation porte sur analyses de programmes Java, avec un coût supplémentaire
de provenance~\cite{provenance}.

\textbf{Proposition.} Envisager Datalog pour plusieurs règles interdépendantes
et une dérivation récursive difficile à maintenir dans le lecteur à règles.
Un graphe de dépendances
de preuves convient à cette question. La récursion seule ne suffit pas à choisir
Datalog puisque SQL en fournit également. Une dérivation explique les règles
appliquées ; elle ne certifie ni les données initiales ni l'intention du propriétaire.

\section{Politiques d'autorisation et états}
\subsection{Formaliser l'attendu avant d'interpréter l'observé}
\textbf{Proposition.} Pour les cas où les droits sont centraux, commencer par
une table lisible : principal, rôle ou attribut, ressource, opération autorisée,
conditions et version de politique. Une fonction déterministe suffit souvent
sur un domaine fini. Ajouter un état uniquement si la propriété le nécessite,
par exemple un rôle révoqué ou une autorisation expirée. Une transition est
alors un changement documentaire d'état, pas une instruction à exécuter sur
un équipement.

Le test reçoit deux couches séparées : politique attendue fournie et trace de
comportement fournie. Un verdict de conformité compare les deux pour la même
identité, la même ressource et le même instant. Il doit rester insuffisant si
l'attribut nécessaire, la portée ou la version sont inconnus. Déduire une
politique attendue uniquement depuis le comportement observé ferait disparaître
les écarts que l'on cherche à mesurer.

\subsection{OPA/Rego : politique déclarative générale}
\textbf{Fait sourcé.} Rego calcule des décisions sur des données structurées.
Sa documentation distingue valeurs non définies, défauts explicites, négation
et quantification universelle. Elle explique qu'une condition satisfaite pour
un élément ne prouve pas qu'elle l'est pour tous~\cite{opa}.

\textbf{Proposition.} OPA convient comme variante de conformité documentaire
si l'équipe souhaite exprimer et versionner des règles indépendamment du code.
Le moteur recevrait uniquement des données locales fournies. La forme proposée
du résultat conserve verdict, motifs, références et obligations manquantes ;
elle exige une règle d'abstention au-dessus des valeurs booléennes.
Une comparaison avec le lecteur à règles doit garder les prédicats métier identiques : traduire
une bonne politique dans Rego puis la comparer à une mauvaise politique Python
ne mesurerait pas la valeur du langage. Le risque principal est le coût de
traduction, de revue et de mise à jour, encore inconnu pour LANCE.

\subsection{Cedar : un modèle centré sur les droits}
\textbf{Fait sourcé.} Cedar structure l'autorisation autour du principal, de
l'opération, de la ressource et du contexte. Les interdictions applicables
l'emportent sur les permissions ; l'absence de permission applicable entraîne
un refus. L'article étudie validation typée et compilation symbolique pour
comparer des politiques. Son évaluation temporelle utilise trois applications
et leurs données en mémoire ; elle exclut certaines étapes d'initialisation
et de récupération~\cite{cedar}.

\textbf{Proposition.} Réserver Cedar à un sous-corpus où une spécification
d'autorisation est disponible et où relations entre entités et exceptions
justifient ce langage. Il peut fournir un évaluateur attendu ou un comparateur
de deux versions de politique. Un refus par défaut du moteur d'autorisation
est une décision d'application ; il ne devient pas une preuve que le dossier
d'audit démontre ce refus dans le déploiement réel. La version du moteur et le
schéma de chaque expérience devront être figés.

\textbf{Choix proposé.} Ne pas démarrer simultanément OPA et Cedar. Tester
d'abord une table de droits et le lecteur à règles proposé. Choisir ensuite un langage si une propriété
concrète manque ou si son coût de maintenance semble acceptable, puis mesurer
ce coût. Une spécification produite à la main constitue un apport humain
qui doit être compté pour tous les comparateurs qui la reçoivent.

\section{Vérification statique de configurations réseau}
Les approches suivantes forment une piste distincte : elles exploitent des
configurations et des propriétés explicites. Elles ne constituent pas des
remplaçants directs de l'interprétation de dossiers incomplets. Leur intérêt
pour LANCE dépend d'abord de la présence d'exports fidèles et d'un domaine de
syntaxe couvert. Une comparaison sur des captures réseau seules serait inadaptée.

\subsection{FIREMAN : anomalies de règles et intention explicite}
\textbf{Fait sourcé.} FIREMAN modélise listes et chaînes de règles, puis
ensembles de paquets avec diagrammes de décision binaires. Il recherche
violations de politique, incohérences et inefficacités. Sa discussion distingue
anomalie formelle et intention de l'administrateur : certains résultats restent
des avertissements. L'évaluation examine des configurations disponibles de
plusieurs types de pare-feu~\cite{fireman}.

\textbf{Proposition.} Utiliser cette famille pour définir un petit ensemble de
propriétés de configuration hors ligne : règle masquée, traitement contradictoire,
permission dépassant une politique fournie. Commencer par les règles et leur
sémantique déterministe, pas par la réimplémentation de l'ensemble du toolkit.
Une redondance volontaire pour défense en profondeur ne doit pas recevoir
automatiquement le label « vulnérabilité ». Le modèle doit préserver l'ordre,
les actions et les règles implicites effectivement retenues.

\subsection{Header Space Analysis : ensembles et transformations}
\textbf{Fait sourcé.} HSA représente les en-têtes comme espaces de bits et
compose des fonctions de transfert, afin d'examiner accessibilité, boucles et
isolation. L'article évalue le réseau Stanford et d'autres configurations.
Sa discussion limite l'analyse aux propriétés du modèle statique : elle
n'explique pas la construction des tables et traite mal les changements plus
rapides que la collecte de snapshots~\cite{hsa}.

\textbf{Proposition.} Une propriété d'isolation sur export complet pourrait
motiver un modèle ensembliste. Sur une simple relation d'adjacence, une fermeture
transitive SQL suffit souvent. Les transformations d'en-tête et la priorité
de règles rendent un modèle de paquets plus pertinent ; l'adjacence seule
ne prouve alors pas l'accessibilité. La propriété reste « dans le modèle
fourni », avec hypothèses et fonctionnalités omises visibles.

\subsection{Batfish et NetPlumber : deux besoins différents}
\textbf{Faits sourcés.} L'article Batfish 2015 décrit un modèle déclaratif
du plan de contrôle, un plan de données dérivé et des propriétés vérifiées avec
solveur. L'évaluation concerne deux réseaux universitaires ; les limites
incluent environnements incomplets, fonctionnalités non modélisées et défauts
matériels ou logiciels hors modèle~\cite{batfish}. La documentation pybatfish
0.36.0 décrit des statuts et avertissements de parsing~\cite{batfishdocs}.
NetPlumber met à jour un graphe de dépendances entre règles après modification ;
ses expériences portent sur trois réseaux. Ses limites incluent états dynamiques
de middleboxes et coût des modifications de liens~\cite{netplumber}.

\textbf{Proposition.} Batfish est une option d'analyse d'exports compatibles,
après contrôle des avertissements. NetPlumber devient pertinent seulement si
une question scientifique porte sur recalcul incrémental ; un audit ponctuel
n'en justifie pas le coût d'intégration. L'article historique et la version
courante d'un logiciel ne sont pas assimilés : le futur manifeste devra figer
versions du serveur, client et fonctionnalités réellement utilisées.

\section{Hiérarchie proposée et critères de sélection}
Le tableau~\ref{tab:choices} distingue couverture et travail d'intégration.
Les priorités sont des appréciations argumentées, sans estimation chiffrée
d'heures ni résultat de benchmark.

\begin{table}[htbp]
\centering\small
\begin{tabularx}{\linewidth}{@{}P{27mm}YYY@{}}
\toprule
Méthode & Pertinence / couverture & Intégration à prévoir & Priorité proposée\\
\midrule
D0, procédure fixe & Contrôles documentaires routiniers, contraste de départ. &
Checklist utile, sorties et échecs conservés. & Secondaire ; ne pas en faire
le seul adversaire du LLM.\\
Lecteur à règles, adaptation proposée de D1 & Obligations connues, contexte,
justification, abstention. & Schéma commun et règles revues sur développement. &
Référence principale hors ligne, à préparer.\\
SQL et relations & Attribution, complétude, déduplication ; récursion simple. &
Clés, contraintes, index, provenance des agrégations. & Ajouter si dossier
volumineux ou relations mal tracées.\\
Datalog / Soufflé & Inférences récursives interdépendantes, dérivations. &
Traduction typée, stratégie de négation, restitution. & Sous-corpus justifiant
la récursion ; pas de gain présumé.\\
Table de droits, puis OPA ou Cedar & Conformité à une politique fournie. &
Formalisation métier, identités et versions ; tests de traduction. & Table
d'abord ; un langage ensuite si nécessaire.\\
FIREMAN / HSA, Batfish & Propriétés de configuration et d'isolation dans un modèle. &
Exports complets, parseurs, sémantique, exclusions. & Piste séparée sous condition
de données et compatibilité.\\
NetPlumber & Vérification incrémentale après changements. & Événements et
dépendances fiables, gestion du modèle. & Hors premier pilote ponctuel.\\
Contraintes symboliques & Équivalence ou propriété sur un domaine explicite. &
Encodage, limites du solveur, contre-exemples. & Après échec motivé d'une
énumération finie ou d'une règle simple.\\
\bottomrule
\end{tabularx}
\caption{Choix de recherche proposés. La couverture désigne les propriétés
du modèle ou du dossier, pas une couverture exhaustive d'un réseau réel.}
\label{tab:choices}
\end{table}

\textbf{Contraintes et symbolique.} Pour un domaine local fini d'identités,
ressources et états, comparer les sorties par énumération peut être plus facile
à relire qu'un encodage SMT. Un solveur devient une option lorsque l'espace
de configurations ou les propriétés dépassent cette énumération. Ses issues
\emph{unknown}, timeout ou erreur doivent rester distinctes d'une preuve.
Un contre-exemple doit être interprété et sa faisabilité vérifiée dans le modèle.
Il ne constitue pas une observation terrain. Ce choix est une inférence de
conception à partir des exigences de formalisation, pas une performance mesurée.

\textbf{Apprentissage classique sans LLM.} Une classification supervisée ou
un score d'anomalie pourrait être ajouté ultérieurement si l'on dispose de
labels indépendants, représentatifs et stables. Pour les propriétés précises
retenues ici, règles et modèles sont le premier choix proposé. L'apprentissage
ajoute séparation des données, réglage de seuils et risque de dépendance aux
artefacts du corpus ; sans jeu justifié, le placer dans le premier pilote
disperserait la recherche. Ce rapport n'effectue pas une revue exhaustive de
l'apprentissage classique et n'en conclut pas l'inutilité.

\section{Quatre expériences hors ligne à préparer}
\textbf{Principe commun.} Les entrées, obligations et références cachées sont
figées avant les résultats. Le jeu de développement sert à corriger parseurs,
règles et prompts ; le test réservé demeure séparé par famille d'origine.
Une variante d'un même générateur ne constitue pas une famille indépendante.
Les comparateurs n'accèdent ni aux labels ni aux dossiers de correction.

\textbf{Porte d'acceptation avant le test réservé.} Sur développement uniquement,
faire relire la couverture du lecteur à règles : chaque critère retenu dispose
d'une règle explicite, d'un contexte applicable et d'un exemple positif, d'un
négatif proche et d'un cas incomplet ou inconnu. Vérifier les justifications,
références aux données, contradictions et obligations manquantes, ainsi que
l'absence d'utilisation des identifiants ou labels cachés. Fixer avant cette
revue les critères d'acceptation et les écarts admissibles sur développement ;
consigner toute lacune et sa correction avec le temps humain actif. Un relecteur
distinct des auteurs accepte la couverture ou demande de reprendre le
développement. Puis figer entrées publiques, règles, schéma, normalisation,
obligations, prompts éventuels, budgets et métriques avant d'ouvrir le test.
Les lacunes découvertes au test restent des résultats ; elles ne sont pas
corrigées pour recalculer silencieusement le même contraste. Une nouvelle
version exige un nouveau gel et un nouveau test indépendant, en conservant
les résultats de la version précédente. Cette porte qualifie une référence
forte pour le périmètre annoncé ; elle ne garantit pas sa performance future.

\textbf{T1 -- Interprétation de dossiers et attribution.} Comparer le lecteur
à règles proposé puis ce même lecteur avec index relationnel ; ajouter un
lecteur LLM recevant les mêmes
observations structurées, si l'on souhaite isoler H4 du protocole existant.
Les cas couvrent soutien, réfutation, données absentes, mauvaise identité,
ressource différente, observations anciennes, contradictions et doublons.
L'oracle distingue propriété réelle connue et soutien accessible au dossier.
Hypothèse : l'index peut améliorer l'attribution ou la relecture ; une égalité
des verdicts est attendue si règles et information sont déjà équivalentes.

\textbf{T2 -- Conformité à une politique d'autorisation.} Sur politiques et
traces synthétiques fournies, comparer fonction déterministe sur table de
droits et un évaluateur OPA \emph{ou} Cedar. Donner au lecteur LLM le même
schéma et la même politique explicite. Inclure permission, refus explicite,
révocation documentée, rôle absent et version indéterminée. L'oracle de politique
est établi par revue indépendante et petit domaine énumérable. Séparer exactitude
de la traduction et exactitude de la décision. Une comparaison entre politique
naturelle pour le LLM et politique formelle pour le lecteur à règles est une étude d'utilité
de bout en bout, avec coûts humains ; elle n'isole pas le seul moteur.

\textbf{T3 -- Propriétés de configuration réseau.} Sur exports synthétiques
ou autorisés et figés, étudier une propriété d'isolation ou de conformité de
règles. La référence simple modélise précisément le sous-ensemble retenu.
Une variante statique plus riche reçoit la même information, ou ses entrées
supplémentaires sont déclarées comme facteur distinct. Contrôler cas valide,
violation locale, fonctionnalités non couvertes, export incomplet et version
ancienne. Aucun verdict de modèle n'est compté comme preuve d'accès réel.
Batfish n'est retenu qu'après vérification de compatibilité des exports et
politique préenregistrée de traitement des avertissements.

\textbf{T4 -- Récursion et maintenance des règles.} Préparer des dossiers où
les dépendances documentaires forment plusieurs niveaux. Comparer SQL récursif
et Datalog à prédicats métier égaux. Tester cycles, doublons, entrée manquante,
contradiction et retrait d'une observation entre deux versions. La question est
la correction et la facilité de restitution/maintenance. Ajouter Datalog sans
cas qui distingue sa valeur ne fournirait qu'une traduction supplémentaire.

\textbf{Mesures proposées.} Exactitude des verdicts et motifs ; précision
et rappel d'étayage ; confirmations abusives ; attribution valide ; couverture
décidée ; résultats insuffisants et conflits ; complétude d'exécution ; temps
écoulé, mémoire, calcul et coût du modèle. Publier les comptes et tous les
essais prévus, y compris erreurs et délais dépassés. Une amélioration de précision
obtenue en s'abstenant davantage doit rester visible.

Consigner séparément heures actives de formalisation, annotation, traduction,
développement, correction et maintenance ; ne pas les inventer depuis des
tokens ou des temps machine. Apparier chaque méthode sur le même dossier,
puis analyser par famille. Le pilote est descriptif ; taille du test, effet
minimal utile et marges des garde-fous se fixent avant le test réservé.
Une incertitude compatible avec gain et perte conduit à un résultat indéterminé.

\section{Discussion, limites et questions pour les relecteurs}
La revue est ciblée et ne couvre pas toutes les approches de vérification
formelle ni tous les outils d'audit. Les expériences originales ne sont pas
reproduites. Les articles anciens apportent des mécanismes et des limites,
pas une preuve de compatibilité des outils actuels avec LANCE. Les performances
rapportées dans des corpus hétérogènes ne permettent aucun classement global.
Une documentation consultée sans version figée devra être remplacée par une
référence de logiciel reproductible avant expérimentation.

Les questions prioritaires pour une contre-revue sont : le lecteur à règles couvre-t-il les
obligations métier sans recevoir les solutions cachées ? L'effet annoncé vient-il
du moteur, d'une meilleure spécification ou de données supplémentaires ? Un
parseur commun masque-t-il des erreurs communes ? La distinction inconnu/faux
survit-elle jusqu'aux mesures ? La politique attendue a-t-elle été inventée
depuis l'observation ? Les familles de test sont-elles suffisamment distinctes ?
Le coût humain d'une formalisation est-il enregistré dans tous les groupes ?

L'oracle ne doit pas être seulement le moteur comparé ou un validateur commun.
Des annotations indépendantes, désaccords conservés et cas finis relus offrent
une référence plus contestable et vérifiable. Un défaut partagé entre modèle,
normalisation et oracle peut néanmoins subsister. Des cas synthétiques sont
adaptés à une erreur ciblée, mais ne donnent pas sa fréquence terrain.

\section{Conclusion}
La priorité scientifique est de valider la couverture du D1 existant et de
préparer son adaptateur hors ligne à règles, avec abstention, contexte et
justification. La porte d'acceptation puis le gel précèdent le test réservé.
Les jointures et règles couvrent déjà beaucoup de problèmes d'attribution.
Datalog, OPA/Cedar et les vérificateurs statiques répondent à des propriétés
différentes ; ils s'ajoutent lorsque les données et les cas le justifient.

Le premier travail testable est T1 sur dossiers communs, puis T2 si des
politiques d'autorisation sont disponibles. T3 exige des exports compatibles
et T4 une difficulté récursive réelle. Une absence de gain du LLM, un gain limité
à la restitution ou une référence classique meilleure restent des conclusions
recevables. À ce stade, seuls la revue et ces protocoles sont produits.

\phantomsection
\addcontentsline{toc}{section}{Références et portée de lecture}
\begin{thebibliography}{10}
\small
\bibitem{sql} SQLite. \emph{SELECT ; The WITH Clause}, documentation officielle
non figée, consultée le 7 octobre 2026.
\href{https://www.sqlite.org/lang_select.html}{SELECT} et
\href{https://www.sqlite.org/lang_with.html}{WITH}.
Lecture : jointures, filtres, NULL, déduplication ; exemples récursifs, cycles.
Pas de comparaison de performance SQL reproduite.
\bibitem{souffle} Projet Soufflé. \emph{Datalog Language}, documentation
officielle non figée, consultée le 7 octobre 2026.
\href{https://souffle-lang.org/datalog-language}{Référence du langage}.
Lecture : relations, types, clauses et négation ; aucune version exécutable
testée. Les détails de provenance proviennent de la référence suivante.
\bibitem{provenance} D. Zhao, P. Subotic et B. Scholz. \emph{Provenance for
Large-scale Datalog}, arXiv:1907.05045v1, juillet 2019.
\href{https://arxiv.org/html/1907.05045v1}{Version primaire consultée}.
Lecture ciblée : sections 3.1--3.4, passages 4.2 et 5 ; sémantique, arbres de preuve,
non-dérivation, expérimentation Doop/DaCapo. Preuves mathématiques non vérifiées
intégralement ; aucune reproduction.
\bibitem{opa} Open Policy Agent. \emph{Policy Language}, documentation
officielle non figée, consultée le 7 octobre 2026.
\href{https://www.openpolicyagent.org/docs/policy-language}{Rego}.
Lecture ciblée : valeurs non définies, négation, quantification, défauts.
Les exemples de documentation ne constituent pas un benchmark LANCE.
\bibitem{cedar} J. W. Cutler et collaborateurs. \emph{Cedar: A New Language
for Expressive, Fast, Safe, and Analyzable Authorization}, version étendue
arXiv:2403.04651v1, 7 mars 2024 ; publication OOPSLA1, avril 2024.
\href{https://arxiv.org/html/2403.04651v1}{Version consultée}.
Lecture ciblée : sections 2, 3.3/3.5, 4.3 et 5 ; sémantique, limites du modèle,
protocole d'évaluation. Preuves Lean et artefact non vérifiés.
\bibitem{fireman} L. Yuan et collaborateurs. \emph{FIREMAN: A Toolkit for
FIREwall Modeling and ANalysis}, IEEE Security and Privacy, 2006, p. 199--213.
\href{https://haochen.org/publications/yuan2006fireman.pdf}{PDF des auteurs}.
Lecture ciblée : sections 2, 4.1, 4.2.2--4.3, 5.1--5.3 ; pages PDF 2--3,
7--11. Modélisation, intention et cadre d'évaluation, sans reproduction.
\bibitem{hsa} P. Kazemian, G. Varghese et N. McKeown. \emph{Header Space
Analysis: Static Checking for Networks}, NSDI 2012, p. 113--126.
\href{https://web.stanford.edu/class/cs244/papers/HeaderSpaceAnalysis.pdf}{PDF universitaire}.
Lecture ciblée : sections 2, accessibilité section 4, début section 7 et section
8 ; pages PDF 2--3, 5--6, 9--10, 13. Instantanés et hypothèses, sans reproduction.
\bibitem{batfish} A. Fogel et collaborateurs. \emph{A General Approach to
Network Configuration Analysis}, NSDI 2015, p. 469--483.
\href{https://ratul.org/papers/nsdi2015-batfish.pdf}{PDF d'auteur}.
Lecture ciblée : sections 3, 6 et 7 ; pages PDF 4--6, 9--13 ; mécanisme,
environnement, restrictions et protocoles. Article historique ; il ne décrit
pas toutes les fonctionnalités d'un logiciel actuel.
\bibitem{batfishdocs} Projet Batfish. \emph{Snapshot Input}, documentation
pybatfish affichée 0.36.0, consultée le 7 octobre 2026.
\href{https://batfish.readthedocs.io/en/latest/notebooks/snapshot.html}{Entrées et parsing}.
Lecture : statuts de fichier et avertissements ; serveur non testé, URL
\code{latest} évolutive. Une version de campagne devra être figée.
\bibitem{netplumber} P. Kazemian et collaborateurs. \emph{Real Time Network
Policy Checking Using Header Space Analysis}, NSDI 2013, p. 99--111.
\href{https://www.usenix.org/system/files/conference/nsdi13/nsdi13-final8.pdf}{PDF USENIX}.
Lecture ciblée : section 3.1, protocole section 6 et discussion section 7 ;
pages PDF 2--3, 9--12. Incidence des mises à jour et limites ; aucune reproduction.
\end{thebibliography}
\end{document}
